\begin{graphillustration}{Примеры: удаление листа сохраняет дерево; лес может иметь изолированную вершину.}{fig:06-tree-forest}
\begin{tikzpicture}[font=\small]
\begin{scope}[shift={(1,0)}]
\node[gi vertex] (r) at (0,0) {$r$};
\node[gi vertex] (a) at (-1,-1) {$a$};
\node[gi vertex] (b) at (1,-1) {$b$};
\node[gi selected] (v) at (1.8,-2) {$v$};
\draw[gi edge] (r) -- (a);
\draw[gi edge] (r) -- (b);
\draw[gi edge] (b) -- (v);
\node[gi note] at (0,-2.7) {Дерево $T$, лист $v$};
\end{scope}
\node[gi note] at (4.3,-1) {$\longrightarrow$};
\begin{scope}[shift={(6.6,0)}]
\node[gi vertex] (r) at (0,0) {$r$};
\node[gi vertex] (a) at (-1,-1) {$a$};
\node[gi vertex] (b) at (1,-1) {$b$};
\draw[gi edge] (r) -- (a);
\draw[gi edge] (r) -- (b);
\node[gi note] at (0,-2.7) {Дерево $T-v$};
\end{scope}
\begin{scope}[shift={(10.5,0)}]
\node[gi vertex] (a) at (0,0) {$a$};
\node[gi vertex] (b) at (1,-1) {$b$};
\node[gi vertex] (c) at (2,0) {$c$};
\node[gi vertex] (d) at (2,-1.8) {$d$};
\draw[gi edge] (a) -- (b);
\draw[gi edge] (b) -- (c);
\node[gi note] at (1,-2.7) {Лес: $c=2$};
\end{scope}
\end{tikzpicture}
\end{graphillustration}
